% GENERATED FILE. Edit canonical lessons, then run npm run generate:books.
% canonical-source-hash: fnv1a64:718ee0572f94b7a6
% canonical-lessons: BN-C338-sonod, BN-C338-tothyo, BN-C338-sakkhyo, BN-C338-melano, BN-C338-khondon-kora

\chapter{Documents and Facts}
\label{ch:bn-documents-and-facts}
\hlchaptermodality{338}

\begin{chapteropening}
\textbf{By the end of this chapter:} \emph{I can check a document against the facts.}

Complete the last lesson of chapter 338: I can check a document against the facts.
\end{chapteropening}

\section[\texorpdfstring{\bn{সনদ} (sônôd)}{সনদ (sônôd)}]{\bn{সনদ} (sônôd) — a certificate; a charter}

\label{lesson:BN-C338-sonod}



\begin{quote}
Before the new one: say the Bengali for to identify, then the Bengali for to uncover.
\end{quote}

\subsection*{You'll want to know: \bn{সনদ}}
\textbf{\bn{সনদ}} — \emph{sônôd} — \textquotedblleft{}a certificate; a charter\textquotedblright{}.

\begin{quote}
\bn{চাকরির} \bn{জন্য} \bn{তোমার} \bn{সনদ} \bn{দরকার।}
\end{quote}

\begin{quote}
\emph{chākrir jônyô tomār sônôd dôrkār.}
\end{quote}

\begin{quote}
You need your certificate for the job.
\end{quote}

\begin{cousinweb}
\textbf{\bn{সনদ}} (\emph{sônôd}) is usually traced to Arabic \emph{sanad}, something to lean on and so a warrant, through Persian. In Bengali it is a paper that vouches for you: a certificate of a degree, a skill or a grant.
\end{cousinweb}

\begin{grammarlens}[title={What you've built}]
One more word for giving reasons and opinions.
\end{grammarlens}

\subsection*{Your turn}
\begin{itemize}
\raggedright
  \item \emph{Say it:} \emph{sônôd}
  \item \emph{Say it:} \emph{sônôd}, once more
  \item \emph{Say it:} \emph{chākrir jônyô tomār sônôd dôrkār.} — You need your certificate for the job.
  \item \emph{Recall:} say the Bengali for to identify, then the Bengali for to uncover, then say \emph{sônôd} again
\end{itemize}

\begin{tcolorbox}[breakable,colback=teal!4,colframe=teal!35!black,title={Before you move on}]
What does \bn{সনদ} mean? (\textquotedblleft{}A certificate; a charter\textquotedblright{}.) Say it once more.
\end{tcolorbox}
\section[\texorpdfstring{\bn{তথ্য} (tôthyô)}{তথ্য (tôthyô)}]{\bn{তথ্য} (tôthyô) — information; facts}

\label{lesson:BN-C338-tothyo}



\begin{quote}
Before the new one: say the Bengali for to uncover, then the Bengali for a certificate.
\end{quote}

\subsection*{You'll want to know: \bn{তথ্য}}
\textbf{\bn{তথ্য}} — \emph{tôthyô} — \textquotedblleft{}information; facts\textquotedblright{}.

\begin{quote}
\bn{এই} \bn{তথ্যটা} \bn{কি} \bn{সত্যি}?
\end{quote}

\begin{quote}
\emph{ei tôthyôṭā ki sôtti?}
\end{quote}

\begin{quote}
Is this information true?
\end{quote}

\begin{cousinweb}
\textbf{\bn{তথ্য}} (\emph{tôthyô}) is Sanskrit \emph{tathya}, 'being really so', truth (Monier-Williams). Bengali uses it for facts and information, a modern use of the word.
\end{cousinweb}

\begin{grammarlens}[title={What you've built}]
One more word for giving reasons and opinions.
\end{grammarlens}

\subsection*{Your turn}
\begin{itemize}
\raggedright
  \item \emph{Say it:} \emph{tôthyô}
  \item \emph{Say it:} \emph{tôthyô}, once more
  \item \emph{Say it:} \emph{ei tôthyôṭā ki sôtti?} — Is this information true?
  \item \emph{Recall:} say the Bengali for to uncover, then the Bengali for a certificate, then say \emph{tôthyô} again
\end{itemize}

\begin{tcolorbox}[breakable,colback=teal!4,colframe=teal!35!black,title={Before you move on}]
What does \bn{তথ্য} mean? (\textquotedblleft{}Information; facts\textquotedblright{}.) Say it once more.
\end{tcolorbox}
\section[\texorpdfstring{\bn{সাক্ষ্য} (sākkhyô)}{সাক্ষ্য (sākkhyô)}]{\bn{সাক্ষ্য} (sākkhyô) — testimony; evidence (given by a witness)}

\label{lesson:BN-C338-sakkhyo}



\begin{quote}
Before the new one: say the Bengali for a certificate, then the Bengali for information.
\end{quote}

\subsection*{You'll want to know: \bn{সাক্ষ্য}}
\textbf{\bn{সাক্ষ্য}} — \emph{sākkhyô} — \textquotedblleft{}testimony; evidence (given by a witness)\textquotedblright{}.

\begin{quote}
\bn{সাক্ষী} \bn{আদালতে} \bn{সাক্ষ্য} \bn{দিল।}
\end{quote}

\begin{quote}
\emph{sākkhī ādālôte sākkhyô dilô.}
\end{quote}

\begin{quote}
The witness gave evidence in court.
\end{quote}

\begin{cousinweb}
\textbf{\bn{সাক্ষ্য}} (\emph{sākkhyô}) is Sanskrit \emph{sākṣya}, testimony or evidence (Monier-Williams), the word for what a \bn{সাক্ষী}, a witness, says. It is given with \bn{দেওয়া}.
\end{cousinweb}

\begin{grammarlens}[title={What you've built}]
One more word for giving reasons and opinions.
\end{grammarlens}

\subsection*{Your turn}
\begin{itemize}
\raggedright
  \item \emph{Say it:} \emph{sākkhyô}
  \item \emph{Say it:} \emph{sākkhyô}, once more
  \item \emph{Say it:} \emph{sākkhī ādālôte sākkhyô dilô.} — The witness gave evidence in court.
  \item \emph{Recall:} say the Bengali for a certificate, then the Bengali for information, then say \emph{sākkhyô} again
\end{itemize}

\begin{tcolorbox}[breakable,colback=teal!4,colframe=teal!35!black,title={Before you move on}]
What does \bn{সাক্ষ্য} mean? (\textquotedblleft{}Testimony; evidence (given by a witness)\textquotedblright{}.) Say it once more.
\end{tcolorbox}
\section[\texorpdfstring{\bn{মেলানো} (melāno)}{মেলানো (melāno)}]{\bn{মেলানো} (melāno) — to match up; to check against}

\label{lesson:BN-C338-melano}



\begin{quote}
Before the new one: say the Bengali for information, then the Bengali for testimony.
\end{quote}

\subsection*{You'll want to know: \bn{মেলানো}}
\textbf{\bn{মেলানো}} — \emph{melāno} — \textquotedblleft{}to match up; to check against\textquotedblright{}.

\begin{quote}
\bn{সনদের} \bn{নামটা} \bn{তোমার} \bn{পরিচয়পত্রের} \bn{সঙ্গে} \bn{মেলাও।}
\end{quote}

\begin{quote}
\emph{sônôder nāmṭā tomār pôrichôypôtrer sôngge melāo.}
\end{quote}

\begin{quote}
Check the name on the certificate against your identity card.
\end{quote}

Past (he/she …): \bn{তিনি} \bn{নামগুলো} \bn{মেলালেন}, \emph{tini nāmgulo melālen}.

\begin{cousinweb}
\textbf{\bn{মেলানো}} (\emph{melāno}) is the causative of \bn{মেলা} 'to meet, to tally', of the family of Sanskrit \emph{milana}, meeting (CDIAL 10134 lists Bengali milan): to make two things meet and see whether they agree. The second thing takes -\bn{এর} \bn{সঙ্গে}.
\end{cousinweb}

\begin{grammarlens}[title={What you've built}]
One more word for giving reasons and opinions.
\end{grammarlens}

\subsection*{Your turn}
\begin{itemize}
\raggedright
  \item \emph{Say it:} \emph{melāno}
  \item \emph{Say it:} \emph{melāno}, once more
  \item \emph{Say it:} \emph{sônôder nāmṭā tomār pôrichôypôtrer sôngge melāo.} — Check the name on the certificate against your identity card.
  \item \emph{Recall:} say the Bengali for information, then the Bengali for testimony, then say \emph{melāno} again
\end{itemize}

\begin{tcolorbox}[breakable,colback=teal!4,colframe=teal!35!black,title={Before you move on}]
What does \bn{মেলানো} mean? (\textquotedblleft{}To match up; to check against\textquotedblright{}.) Say it once more.
\end{tcolorbox}
\section[\texorpdfstring{\bn{খণ্ডন} \bn{করা} (khônḍôn kôrā)}{খণ্ডন করা (khônḍôn kôrā)}]{\bn{খণ্ডন} \bn{করা} (khônḍôn kôrā) — to refute}

\label{lesson:BN-C338-khondon-kora}



\begin{quote}
Before the new one: say the Bengali for testimony, then the Bengali for to check against.
\end{quote}

\subsection*{You'll want to know: \bn{খণ্ডন} \bn{করা}}
\textbf{\bn{খণ্ডন} \bn{করা}} — \emph{khônḍôn kôrā} — \textquotedblleft{}to refute\textquotedblright{}.

\begin{quote}
\bn{মীরা} \bn{প্রমাণ} \bn{দিয়ে} \bn{অভিযোগটা} \bn{খণ্ডন} \bn{করল।}
\end{quote}

\begin{quote}
\emph{mīrā prômān diye ôbhijogṭā khônḍôn kôrlô.}
\end{quote}

\begin{quote}
Mira refuted the charge with evidence.
\end{quote}

Past (he/she …): \bn{তিনি} \bn{যুক্তিটা} \bn{খণ্ডন} \bn{করলেন}, \emph{tini juktiṭā khônḍôn kôrlen}.

\begin{cousinweb}
\textbf{\bn{খণ্ডন} \bn{করা}} (\emph{khônḍôn kôrā}) is Sanskrit \emph{khaṇḍana}, breaking into pieces and, in argument, refuting (Monier-Williams), with \bn{করা}: to break a claim apart.
\end{cousinweb}

\begin{grammarlens}[title={What you've built}]
One more word for giving reasons and opinions.
\end{grammarlens}

\subsection*{Your turn}
\begin{itemize}
\raggedright
  \item \emph{Say it:} \emph{khônḍôn kôrā}
  \item \emph{Say it:} \emph{khônḍôn kôrā}, once more
  \item \emph{Say it:} \emph{mīrā prômān diye ôbhijogṭā khônḍôn kôrlô.} — Mira refuted the charge with evidence.
  \item \emph{Recall:} say the Bengali for testimony, then the Bengali for to check against, then say \emph{khônḍôn kôrā} again
\end{itemize}

\begin{tcolorbox}[breakable,colback=teal!4,colframe=teal!35!black,title={Before you move on}]
What does \bn{খণ্ডন} \bn{করা} mean? (\textquotedblleft{}To refute\textquotedblright{}.) Say it once more.
\end{tcolorbox}
